\subsection{KV Master: Request-Specific Backup Decisions}
\label{sec:kv-placement}
\label{sec:kv-master}

KV Master (3) makes checkpoint decisions, choosing recovery tiers, each request's
Peer depth $L_r$, and historical hot blocks. The Recovery Manager fixes
the retry/Peer target at arrival; the target Adapter publishes only a
window-level Peer capacity budget (\S\ref{sec:kv-runtime}).
Different requests can use different depths within the same window.

\paragraph{Decision triggers and checkpoint interval.}
The Master establishes a policy at arrival, offers sealed state at
prefill completion, and reviews an active request after each additional
$Z$ generated tokens. The block-aligned interval $Z$ is a fixed system
configuration, not an adaptive output. Significant slack or DAG changes, revised budgets,
and material changes in context or copy conditions can refresh policy
earlier. Sealing updates block eligibility without triggering a depth
decision for every token. Completion releases that request's Peer
references while preserving shared references; unchanged decisions
produce no new transfer. Hot-block statistics use the preceding 512
decoding steps, independently of checkpoint and capacity-control windows.

A trigger requests a decision, not a completed replica. The replay tail
contains known tokens beyond the latest valid committed frontier;
queued copies can make it longer than $Z$. Tier R need not checkpoint.
The Master uses reported commit progress when estimating this tail.

\paragraph{Choose protection tiers.}
Let $T_{\mathrm{rc}}$ be historical-prefix recomputation cost,
$T_{\mathrm{rp}}$ known-tail replay cost, and $W_0$ exposed DRAM-loading
wait without Peer residency. Detection, rerouting, and readmission
contribute $R_{\mathrm{floor}}$. The Master uses
\begin{equation}
\begin{aligned}
 R_{\mathrm{base}}&=R_{\mathrm{floor}}+T_{\mathrm{rp}},\\
 R_R&=R_{\mathrm{base}}+T_{\mathrm{rc}},\quad
 R_D=R_{\mathrm{base}}+W_0.
\end{aligned}
\label{eq:recovery-base}
\end{equation}
Costs are profiled jointly by model, committed context, replay tail,
hardware, and representative concurrency.
Tier \emph{R} requires no backup when $\sigma_r\ge R_R$;
tier \emph{D} retains DRAM when $R_D\le\sigma_r<R_R$;
tier \emph{P} retains DRAM and requests Peer protection otherwise.
Tier changes use 20\% hysteresis around these nominal boundaries.
A tier-P request without a Peer grant retains DRAM recovery.
Slack orders competing requests; it does not scale their hot-set
fraction or impose a serial retry order.

\paragraph{Define and rank protection objects.}
A protection object is a \emph{(layer, block)} containing all KV heads;
the profiled selectors use 64-token blocks. Sealing makes a full block
eligible for copying; the commit frontier determines the prefix that
survives failure. Known uncommitted tokens are replayed without resampling.

The historical blocks needed by a length-$n$ replay form the union
$U_l(n)=\bigcup_{q=1}^{n}\bigcup_g S_{l,q,g}$, where $S_{l,q,g}$ is the
native selector output. Repeated reads count once. A fixed preload
fraction $h$ is configured for each model and sparse-attention method;
it is not the selector's instantaneous Top-$k$ fraction or a guaranteed
hit rate. With $n_r(\tau)$ sealed historical blocks per layer, the Master uses
\begin{equation}
 K_r(\tau)=\lceil h n_r(\tau)\rceil
 \label{eq:hot-block-count}
\end{equation}
slots in each layer, including mandatory sealed sink/window blocks.
The configured quota must accommodate them. Remaining slots follow
observed selection counts, breaking ties by recency. An optional block
leaves the candidate set after two successive decoding-statistics windows
outside the top set. A replacement needs fresh admission if the old
replica still has a request reference.

\paragraph{Select a depth for each request.}
At a policy decision, the Master reviews $L_r$ using that request's
committed history, generation progress, slack, and recovery profiles.
Both transfer bytes and replay computation depend on context; hot-set
misses can also vary with the replay tail. Candidate costs average over
plausible future checkpoint phases and commit lag, using the same
request-growth and fault-state samples across depths. This does not
assume that a future failure occurs at the current tail or at $N=Z$.

For each request in slack order, the Master compares a small configured
set of depths within the model, including zero and its current depth. Each candidate
previews incremental admission against the same target-state snapshot,
accounting for shared copies, live references, in-flight reservations,
and copy limits. Recovery estimates credit only expected committed
coverage, including retained old copies. The score adds predicted
foreground delay from the candidate's backup work and memory occupancy
to recovery delay beyond the request's slack, weighted by source-failure
probability from the configured risk model over the same policy horizon. It is a local decision
heuristic under slack priority, not a global optimum or workflow-makespan sum.

The minimum-score candidate wins, with smaller depths breaking ties.
An existing depth changes only when its advantage exceeds a configured
uncertainty margin. Missing applicable profiles retain an existing target;
a new request starts with zero Peer depth and DRAM protection.
Only the chosen plan updates the preview state before the next request.
Real allocation rechecks availability atomically; it may admit fewer
objects than the preview.

For uniform layer layouts, a tier-P request offers the first $L_r$
hot sets: $L_rK_r(\tau)$ logical layer-block objects plus required state.
Tiers R/D have a zero target depth. The target is held between policy
decisions, while the hot-block quota grows with sealed history.
D/P keep per-layer hot sets as recovery-preload hints. Actual Peer
coverage can be partial; lowering $L_r$ does not revoke references to
old copies (\S\ref{sec:kv-runtime}).

\begin{figure}[t]
  \raggedright
  \refstepcounter{designalgorithm}\label{alg:placement}
  \noindent\textbf{Algorithm \thedesignalgorithm: Per-request backup policy}
  \par\smallskip
  \noindent Input: requests, slack, profiles, targets, budgets, copy state.
  Output: DRAM set $D_s$, depths $\{L_r\}$, and Peer plans $\{Q_p\}$.
  \par\smallskip
  \begin{tabularx}{\columnwidth}{@{}r@{\quad}X@{}}
    \toprule
    1 & $D_s,H\gets\varnothing$; $Q_p\gets[\,]$ for each target $p$ \\
    2 & Snapshot retained copies and reserved capacity as $S$ \\
    3 & \textbf{for} $r$ in $(\sigma_r,\mathrm{ID}_r)$ order: \\
    4 & \quad$\tau_r\gets\mathrm{Tier}(r,R_R,R_D,\sigma_r)$ \\
    5 & \quad\textbf{if} $\tau_r=R$: $L_r\gets0$; \textbf{continue} \\
    6 & \quad Add $\mathrm{sealed}(r)\cup\mathrm{state}(r)$ to $D_s$ \\
    7 & \quad$H_{r,l}\gets\mathrm{HotSet}(r,l,K_r)$ for each layer \\
    8 & \quad\textbf{if} $\tau_r=D$: $L_r\gets0$; \textbf{continue} \\
    9 & \quad\textbf{for each} candidate depth $d$: \\
    10 & \qquad$G_d\gets\mathrm{Preview}(S,r,H_{r,1:d})$ \\
    11 & \qquad Score $G_d$; keep missing costs unknown \\
    12 & \quad$L_r\gets\mathrm{StableBest}(r,\{G_d\})$ \\
    13 & \quad Append $(r,L_r,H_{r,1:L_r},G_{L_r})$ to $Q_{p_r}$ \\
    14 & \quad Update $S$ with the chosen plan only \\
    15 & \textbf{return} $D_s,H,\{L_r\},\{Q_p\}$ \\
    \bottomrule
  \end{tabularx}
  \Description{KV Master chooses recovery tiers, hot blocks, and a
    separate Peer depth for each request. Depth previews share a common
    snapshot within a request; only its selected plan advances the
    snapshot. Target Adapters enforce capacity and report committed copies.}
\end{figure}

Algorithm~\ref{alg:placement} shows a full refresh; a token-interval
event updates only the affected request while retaining other policies
and reservations. \textsc{Tier} applies hysteresis, and
\textsc{StableBest} implements the switching margin and profile fallback.
The state snapshot is virtual, not a physical reservation.
Sorting costs $O(n\log n)$, followed by bounded candidate scans.
Runtime allocation deduplicates shared objects and reconciles policy
with committed copies, reporting actual residency and commit progress.

\paragraph{Commit and recover by actual readiness.}
Peer and fault-independent DRAM paths commit separately, after transfer
completion and model/layout/shard validation. Low-priority copying is
rate-limited (default 20\% of link capacity), with target HBM bandwidth
also accounted for. Recovery reads committed objects, preloads hot
blocks in layer order, and gives newly revealed misses priority over
speculative copies. Selector state arrives before selection; each layer
waits for all required inputs. Actual DMA granularity can limit copy
preemption. Tier R rebuilds history. Placement changes neither selector
outputs nor known replay tokens and adds no attention approximation;
compatible weights, layout, and operators reconstruct the same logical
KV state, without promising identical future stochastic samples across
floating-point schedules.

\begin{figure}[t]
  \centering
  \includegraphics[width=\columnwidth]{figures/design-single/design-layer-pipeline.pdf}
  \caption{Modeled recovery pipelines: MiniCPM4-8B, 32 layers, 32K
    history, 64-token known tail; $C=2.87$\,ms/layer (measured) and
    $D=5.37$\,ms at an assumed 6.25\,GB/s. In (d), 147\,MiB of Peer
    residency removes preload wait. Totals end at tail-KV reconstruction.
    The homogeneous, preemptible-copy model excludes detection,
    queuing, and interference.}
  \label{fig:design-pipeline}
  \Description{Four timelines compare token-block transfer, full-layer
    pipelining, hot preloads with demand fetch, and Peer-resident hot sets.}
\end{figure}

Figure~\ref{fig:design-pipeline} illustrates the depth tradeoff at one
fixed tail length: layer pipelining, hot preloads, and seven resident
hot sets reduce exposed wait from 172 to 83, 44, and 20\,ms.
This calibrated example does not validate request-specific depth selection
or make seven layers appropriate for every request and checkpoint state.
